\documentclass[11pt,a4paper]{article}

\usepackage[margin=2.3cm]{geometry}
\usepackage{booktabs}
\usepackage{amsmath}
\usepackage{xcolor}
\usepackage{caption}
\usepackage{tabularx}
\usepackage{enumitem}
\usepackage[hidelinks]{hyperref}

\captionsetup{font=small,labelfont=bf}
\definecolor{accent}{RGB}{140,21,21}
\newcommand{\code}[1]{\texttt{\small #1}}
\newcommand{\hi}[1]{\textbf{\textcolor{accent}{#1}}}

\title{\vspace{-1.4cm}\textbf{Why This Study Evaluates Chronologically}\\[0.3cm]
\large The argument, the published support, and our own controlled measurement\\[0.4cm]}
\author{Divyaraj Chudasama \\ Indian Institute of Technology Hyderabad}
\date{\today}

\begin{document}
\maketitle

\section{The claim}
This study evaluates on a \emph{chronological} split --- train on the earliest transactions,
test on the latest --- rather than the random stratified split used by most published work on
the ULB dataset.

The justification is a \hi{validity} argument, not a performance one. A chronological split is
correct because it is the only protocol whose question has a deployment counterpart. Whether it
also produces lower scores is a separate, empirical matter, addressed in
Section~\ref{sec:evidence}.

\section{Why a random split asks the wrong question}
A random split assigns transactions to train and test by coin flip. The model is therefore
fitted on transactions that occurred \emph{after} the ones it is scored on. No deployed fraud
system can do this: at scoring time, the future has not happened yet. Three consequences are
measurable directly on ULB.

\paragraph{(a) The distribution moves, even inside 48 hours.}
ULB spans exactly $48$ hours. Across our chronological partitions the fraud prior falls by
$37\%$ relative and the mean transaction amount by $19\%$:

\begin{center}\small
\begin{tabular}{lcccc}
      \toprule
      Partition & Window (h) & Frauds & Fraud rate & Mean amount\\
      \midrule
      Train & $0.00$--$36.92$ & 384 & $0.193\%$ & $89.77$\\
      Validation & $36.92$--$42.04$ & 56 & $0.131\%$ & $97.52$\\
      Test & $42.04$--$48.00$ & 52 & $0.122\%$ & $72.54$\\
      \bottomrule
\end{tabular}
\end{center}

\noindent The window also covers two full day--night cycles with a $35.5\times$ ratio between
the peak and trough hourly fraud rates ($02$:$00$ peaks at $1.71\%$). Shuffling destroys all of
this structure by construction: it makes the train and test distributions identical by
definition, which is precisely the condition a deployed model does not enjoy.

\paragraph{(b) Shuffling separates duplicate transactions into train and test.}
ULB contains $1{,}081$ exact duplicate rows. All $773$ duplicate groups share a single
\code{Time} value, so chronological sorting places every copy adjacent to its twin and they
always land in the same partition. Measured across all three boundaries:

\begin{center}\small
\begin{tabular}{lcccc}
      \toprule
      & Train$\to$Test & Train$\to$Val & Val$\to$Test & Total leaked\\
      \midrule
      Chronological split & 0 & 0 & 0 & \hi{0}\\
      Random split & 249 & 207 & 61 & \hi{517}\\
      \bottomrule
\end{tabular}
\end{center}

\noindent Under shuffling the model is fitted on a row and then scored on a byte-identical copy
of it. That is memorisation being reported as generalisation, and the chronological protocol is
immune to it for free --- without deleting any data.

\paragraph{(c) The resulting estimate has no operational meaning.}
A random-split score answers ``how well would this model rank transactions if it had already
seen the future?'' That number cannot be used to size a review queue, set an alert threshold, or
forecast caseload, because the condition it assumes never holds in production.

\section{What the published literature says}
The four sources below support the \emph{principle}. We deliberately do not compare our absolute
scores to theirs --- see Section~\ref{sec:nomatch}.

\begin{itemize}[itemsep=4pt]
      \item \textbf{Dal Pozzolo, Boracchi, Caelen, Alippi and Bontempi (2018)}, \emph{Credit Card
            Fraud Detection: A Realistic Modeling and a Novel Learning Strategy}, IEEE
            Transactions on Neural Networks and Learning Systems 29(8):3784--3797. \\
            \small\url{https://re.public.polimi.it/bitstream/11311/1044896/1/08038008.pdf}\\
            \normalsize These are the authors of the ULB dataset itself. They formalise fraud
            detection around three operating conditions --- \emph{concept drift} (customer habits
            evolve and fraudsters change strategy over time), class imbalance, and
            \emph{verification latency} (only a small set of transactions is checked in time) ---
            and argue that a realistic evaluation must reflect them. Concept drift and
            verification latency are both inherently temporal; a shuffled split represents
            neither.

      \item \textbf{Sang (2026)}, \emph{Enhancing Credit Card Fraud Detection under Severe Class
            Imbalance using Cost-Sensitive Learning and Threshold Optimization}, EAI Endorsed
            Transactions on ISMLA, Vol.~3. \\
            \small\url{https://publications.eai.eu/index.php/ismla/article/view/12078}\\
            \normalsize The only study we located that evaluates ULB under \emph{both} protocols.
            It reports that the chronological split yields lower PR-AUC than the random one and
            concludes that ``random stratification may provide optimistic estimates for this
            benchmark dataset when earlier and later transactions are mixed across training,
            validation, and test partitions.'' It also lists, among the gaps motivating the work,
            that ``random stratified validation can mix earlier and later transactions,
            potentially overstating future-period performance,'' and describes its own
            chronological experiment as evaluating ``deployment realism.''

      \item \textbf{\emph{Data Leakage and Deceptive Performance: A Critical Examination of
            Credit Card Fraud Detection Methodologies}} (2025), arXiv:2506.02703. \\
            \small\url{https://arxiv.org/abs/2506.02703}\\
            \normalsize An audit of published ULB pipelines. It identifies four recurring flaws,
            of which the third is \emph{``inadequate temporal validation for transaction
            data.''} It demonstrates the stakes by constructing a deliberately leaky minimal MLP
            that reaches $99.9\%$ recall and outperforms far more sophisticated published
            methods, concluding that ``proper evaluation methodology matters more than model
            complexity in fraud detection research.''

      \item \textbf{Khaleghpour and McKinney}, \emph{Leakage Safe Graph Features for
            Interpretable Fraud Detection in Temporal Transaction Networks}, arXiv:2603.06632. \\
            \small\url{https://arxiv.org/abs/2603.06632}\\
            \normalsize Applies a strict chronological split to the Elliptic dataset --- a
            different dataset, cited here only for the principle --- on the grounds that it
            ``reflects a realistic deployment setting in which models are trained on historical
            data and evaluated on future transactions.'' It observes that the illicit rate
            decreases across time, ``indicating temporal distribution shift that can affect
            generalization,'' and reports a substantial drop from validation to the held-out
            future test period.
\end{itemize}

\section{Our own controlled measurement}\label{sec:evidence}
The literature establishes the principle. To measure the \emph{size} of the effect we compare
our own models against themselves, holding everything fixed --- same data, same features, same
code, same hyperparameters, same threshold rule --- and changing only how rows are assigned to
partitions. This is the only comparison in which the split is the sole variable.

\begin{table}[h]
      \centering\small
      \caption{Chronological versus random, everything else held constant.}
      \label{tab:controlled}
      \begin{tabularx}{\textwidth}{Xccc}
            \toprule
            Configuration & chronological & random & random $-$ chrono\\
            \midrule
            \multicolumn{4}{l}{\emph{at $70/15/15$}}\\
            XGBoost + class weight & 0.7688 & 0.8372 & $+0.0684$\\
            XGBoost + SMOTE-ENN & 0.7766 & 0.8104 & $+0.0338$\\
            CatBoost + GAN & 0.7804 & 0.8399 & $+0.0595$\\
            MLP raw & 0.7794 & 0.8267 & $+0.0473$\\
            LSTM + class weight & 0.7803 & 0.7965 & $+0.0162$\\
            RF + SMOTE-ENN & 0.7947 & 0.8120 & $+0.0173$\\
            GraphSAGE + weight & 0.7389 & 0.6948 & $-0.0441$\\
            \cmidrule(l){2-4}
            & \multicolumn{3}{r}{\emph{mean $+0.0283$, positive in 6 of 7}}\\
            \midrule
            \multicolumn{4}{l}{\emph{at $66/14/20$}}\\
            XGBoost + class weight & 0.8055 & 0.8291 & $+0.0236$\\
            XGBoost + SMOTE-ENN & 0.8102 & 0.8020 & $-0.0082$\\
            RF + SMOTE-ENN & 0.8224 & 0.7907 & $-0.0317$\\
            \cmidrule(l){2-4}
            & \multicolumn{3}{r}{\emph{mean $-0.0054$, positive in 1 of 3}}\\
            \bottomrule
      \end{tabularx}
\end{table}

Pooling all ten paired measurements:
\[
  \text{random} - \text{chronological} \;=\; +0.0182,\qquad
  \text{95\% CI } [-0.0084,\, +0.0448],\qquad
  \text{positive in } 7/10 .
\]

\paragraph{Honest reading.} The direction agrees with theory and with every source in Section~3,
and the point estimate is positive. But the interval includes zero ($t = 1.55$, one-sided sign
test $p = 0.172$), and the estimate is not stable across cut points. \hi{Our data show the
direction, not a reliable magnitude.} Two further notes bound the interpretation:

\begin{itemize}[itemsep=2pt]
      \item Our random arm carries the $517$ leaked duplicate rows documented above, which
            inflates it. The measurement is therefore biased \emph{towards} finding a larger
            effect, not a smaller one.
      \item Our random arm applies all resampling \emph{inside} the fitted pipeline, so it
            captures the ordering effect only. It excludes the resample-before-split channel that
            arXiv:2506.02703 identifies as common in published work, which is the larger of the
            two.
\end{itemize}

\noindent Neither of these changes the conclusion of Section~1: the case for chronological
evaluation rests on validity, and would stand even if the two protocols scored identically.

\section{Why we do not match numbers with the published studies}\label{sec:nomatch}
Absolute PR-AUC figures are not comparable across these papers, and this study does not treat
them as though they were. The reasons are concrete:

\begin{itemize}[itemsep=3pt]
      \item \textbf{Different architectures and objectives.} Sang (2026) uses cost-sensitive
            XGBoost and an MLP trained with weighted BCE plus Focal Loss, with synthetic
            oversampling deliberately excluded. This study's champion is a Random Forest with
            SMOTE-ENN. These are different models solving the problem in different ways.
      \item \textbf{Different search budgets and spaces.} Sang tunes with Optuna over tree depth,
            minimum child weight, learning rate, subsampling ratios, regularisation coefficients
            and the positive-class weight; the trial count is not stated. This study uses 30
            Optuna trials by default, and on the random arms varied only that one weighting
            parameter over a six-point grid. A difference in score between the two is at least as
            likely to reflect search effort as model quality.
      \item \textbf{Undocumented preprocessing.} Several studies do not state where resampling
            sits relative to the split, or whether duplicates were removed. Those choices move
            the metric more than the model does, and cannot be reconstructed from the papers.
      \item \textbf{Different data.} Some headline figures are obtained on private datasets with
            far more fraud cases, where interval widths are an order of magnitude tighter for
            reasons of sample size alone.
      \item \textbf{Measurement precision.} On this test set a single transaction is worth
            $\pm0.020$ PR-AUC and bootstrap intervals are $0.196$--$0.236$ wide. Differences of
            the size usually discussed between papers are smaller than the error bars on either
            side of the comparison.
\end{itemize}

\noindent For these reasons the evidence offered here is Table~\ref{tab:controlled}, in which
the split is the only thing that changes, rather than any comparison of this study's absolute
scores against a published leaderboard.

\section{Conclusion}
\begin{enumerate}[itemsep=3pt]
      \item Chronological evaluation is used because it is the only protocol whose question a
            deployed system could answer. This holds regardless of the scores it produces.
      \item On ULB specifically, shuffling destroys measurable temporal structure and separates
            $517$ duplicate rows across the train/test boundary; the chronological split leaks
            none.
      \item The dataset's own authors, and three subsequent studies, all argue for temporally
            faithful evaluation on the same grounds.
      \item Our controlled measurement puts the effect at $+0.018$ PR-AUC with a confidence
            interval that includes zero. The direction is consistent; the magnitude is not
            established at this sample size, and we do not claim it is.
\end{enumerate}

\end{document}
